\documentclass[reprint,onecolumn,groupedaddress,amsmath,amssymb]{revtex4-2}

\usepackage[]{hyperref}
\usepackage{graphicx}
\usepackage{enumitem}
\usepackage{xr}

\hypersetup{
  colorlinks   = true,
  urlcolor     = blue,
  linkcolor    = blue,
  citecolor   = red
}

\newcommand{\Nc}{N_\mathrm{c}}
\newcommand{\Nb}{N_\mathrm{B}}

\begin{document}

\title{Equilibria with chemical reactions}

\author{David Zwicker}
\affiliation{Max Planck Institute for Dynamics and Self-Organization, Am Fa{\ss}berg 17, 37077 G{\"o}ttingen, Germany}

\maketitle

We consider a system of $\Nc$ components, $X_1, \ldots, X_{\Nc}$, that can undergo $\Nb$ different reactions,
\begin{equation}
    \sum_{i=1}^{\Nc} \sigma^{(b)}_{\rightarrow, i} X_i
    \rightleftharpoons
    \sum_{i=1}^{\Nc} \sigma^{(b)}_{\leftarrow, i} X_i
\end{equation}
for each reaction $b=1,\ldots, \Nb$.
The stoichiometric coefficients count molecules, as is standard for chemical reaction
equations. Writing $\sigma^{(b)}_i =
\sigma^{(b)}_{\leftarrow, i}-\sigma^{(b)}_{\rightarrow, i}$, equilibrium requires
the corresponding molecular chemical potentials to balance. In the package, $\mu_i$
denotes the chemical potential per volume, so the molecular chemical potential is
$l_i\mu_i$, where $l_i=\nu_i/\nu$ is the relative molecule volume:
\begin{equation}
    \sum_{i=1}^{\Nc} \sigma^{(b)}_i l_i \mu_i = 0
\end{equation}

Volume fractions are proportional to molecule numbers times molecule volumes.
Consequently, a reaction extent changes the volume fractions according to
\begin{equation}
    \Delta \phi_i = l_i \sigma^{(b)}_i \Delta \xi_b \;.
\end{equation}
Thus, the matrix used to find conserved volume-fraction combinations is the
stoichiometric matrix with each component column weighted by $l_i$. This weighting
converts molecule-number coefficients into volume-fraction changes; it does not change
the molecular-number convention of the reaction equation. In particular, a reaction
conserves volume when $\sum_i l_i \sigma^{(b)}_i=0$.

Each reaction (in general) leaves conserved linear combinations of volume fractions,
\begin{equation}
    \sum_{i=1}^{\Nc} q^{(b)}_i \phi_i = \psi_b
\end{equation}
where $q^{(b)}_i$ are linearly independent vectors in the null space of the size-weighted
stoichiometric matrix,
\begin{equation}
    \sum_{i=1}^{\Nc} q^{(b)}_i l_i \sigma^{(b)}_i = 0
    \;.
\end{equation}
Once we determine a basis for this null space, e.g., using
\texttt{scipy.linalg.null\_space}, we can use it to express a linear constraint.
Moreover, we need to impose a grandcanonical potential for the respective species.
The imposed chemical potential difference relates to the internal energy difference
of the species, which is not part of the free energy in our formulation.


\end{document}